\section{Introduction}
\label{sec:intro}

Deciding whether a research idea is new is not a property of the idea alone. The
same proposal is a contribution in one year and a restatement in the next, and
the difference lies entirely in what has been published in between. Reviewers,
program committees, and, increasingly, automated idea-generation systems all
make this judgment, and they make it by reading the idea against the record of
prior work. As machine-generated research proposals become cheap to produce,
the ability to assess them at scale has become a bottleneck: a system that
proposes a thousand ideas is only useful if something can tell which of them
are worth pursuing. This motivates a task we state explicitly as
\emph{evidence-conditioned} novelty assessment, in which a model is asked to
grade a target idea relative to a given, temporally admissible set of prior
work, rather than against an unspecified notion of what the field knows.

% TODO: replace the schematic A/B example below with the verified real case
% once the case study in Section~\ref{sec:case} is finalized.
Most existing approaches reduce this judgment to similarity: retrieve the
nearest prior papers and score how close the target is to the closest of them.
This misses how research contributions are actually distributed. The pieces of
a target idea are typically spread across several papers rather than
concentrated in one. A mechanism may already appear in one line of work, the
training signal it is paired with in another, and the evaluation setting in a
third, while no single paper is especially similar to the target as a whole.
Under a nearest-neighbor comparison such an idea looks novel, because the
maximum over single-document similarities is low; under any reading of the
literature it is largely covered. What is needed is a notion of \emph{joint}
coverage, in which several prior works together account for different parts of
the same target.

Joint coverage alone, however, is still a scalar, and a scalar cannot say
\emph{which} part remains. Two ideas that are covered to the same degree are
not equally novel if one leaves only an implementation detail uncovered while
the other leaves its central mechanism unaccounted for. A second failure is
subtler. Every component of an idea may be individually well covered, and the
idea may nonetheless be new, because the way the components are combined has no
precedent; conversely, a combination may look new simply because no single
retrieved paper contains both parts. Aggregate similarity collapses both
distinctions. This suggests decomposing the target instead of scoring it as a
whole: identify its constituent semantic content, determine which parts the
retrieved evidence supports, and treat what is left over, together with the
interactions among the parts, as the object to be evaluated.

We instantiate this idea as Evidence-Grounded Residual Decomposition (EGRD).
Given a target idea and a set of temporally admissible prior works, EGRD first
models historical coverage, scoring each target--prior pair, ranking candidates
by coverage strength, and aggregating the retained evidence under a
multiple-instance formulation that yields a joint coverage estimate. It then
decomposes the target into latent semantic units through a set of learned
queries, so that each unit attends to a distinct part of the target's content.
The central step aligns these units with the evidence: EGRD learns a support
matrix whose entries measure how strongly each retained prior supports each
unit, combines the per-prior supports into a unit-level coverage through a
noisy-OR aggregation, and defines a \emph{unit residual} as the part of each
active unit that no prior accounts for. A second residual is defined over pairs
of units, using the extent to which a single prior supports both units as an
approximation of whether their combination has precedent, so that novel
interactions among individually known components remain visible. Unit
residuals, interaction residuals, and the coverage context are then fused to
predict a novelty grade. Finally, because the model is trained under a fixed
evidence budget, we introduce an inference-time scheme that partitions a larger
evidence set into blocks, encodes each block independently, and combines the
per-block predictions conservatively, allowing the same trained model to
consume several times more evidence without retraining.

Our contributions are as follows.
\begin{itemize}
  \item \textbf{An evidence-conditioned residual formulation.} We recast
  novelty assessment from measuring overall similarity to locating the joint
  coverage provided by prior work and then pricing what remains, and we give a
  task definition that is explicitly relative to a temporally admissible
  evidence set.
  \item \textbf{A multi-domain, evidence-paired dataset.} We release a dataset
  spanning three research areas, annotated at both the idea level and the
  target--prior pair level under temporal admissibility constraints. Its
  pairing is what makes evidence-conditioned evaluation and evidence
  intervention experiments possible; its labels are model-generated, and we
  describe it as a dataset rather than a benchmark for that reason.
  \item \textbf{Unit and interaction residual modeling.} We design a mechanism
  that aligns latent semantic units with historical support and jointly
  represents the uncovered content and the potentially novel interactions among
  covered components.
  \item \textbf{Budget-scalable conservative evidence aggregation.} We propose
  an inference-time procedure that lets a model trained with a small evidence
  budget consume substantially larger evidence sets by blockwise encoding and
  conservative cross-block aggregation, without retraining.
\end{itemize}

% TODO: add a closing sentence summarizing the frozen empirical results once the
% main table is final. Do not pre-write claims of significance or of
% across-the-board improvement over baselines.
